% generated by code/revisao (Stata 17) -- do not edit by hand
\begin{table}[htbp]\centering\footnotesize
\providecommand{\sym}[1]{\ensuremath{^{#1}}}
\caption{Table S15. Bayesian model averaging over the control set (T3 sample)}\label{tab:TabS_bma}
\begin{tabular}{l*{2}{c}}
\toprule
 & Posterior inclusion probability & Posterior mean \\
\midrule
Log GDP per capita, t\_p $-$ 1 & 1.000 & $-$5.549 \\
Openness, t\_p $-$ 1 & 1.000 & 0.030 \\
Population growth, t\_p $-$ 1 & 0.998 & $-$0.633 \\
Inflation, t\_p $-$ 1 & 0.915 & $-$0.025 \\
Terms-of-trade growth, t\_p $-$ 1 & 0.914 & 0.196 \\
Wage bill (\% of GDP), t\_p $-$ 1 & 0.774 & $-$0.116 \\
Investment, t\_p $-$ 1 & 0.505 & $-$0.027 \\
Human capital, t\_p $-$ 1 & 0.267 & 0.508 \\
Meritocratic recruitment, t\_p $-$ 1 & 0.260 & 0.131 \\
Public debt, t\_p $-$ 1 & 0.115 & $-$0.001 \\
\bottomrule
\end{tabular}
\begin{minipage}{0.95\linewidth}\vspace{4pt}\scriptsize \textit{Notes:} Bayesian model averaging (De Luca and Magnus, 2011) over all 2^10 combinations of the regressors, all treated as auxiliary, after removing country and period fixed effects from every variable (Frisch-Waugh) in the effective T3 sample (346 observations, singletons excluded). The posterior inclusion probability measures robustness to the choice of controls; it says nothing about endogeneity. Weighted-average least squares (WALS) on the same data: wage-bill coefficient -0.131 (t = -2.57).\end{minipage}
\end{table}
